% Generated from the matching Typst scaffold by
% scripts/migrate_typst_report_to_latex.py (prose drafted 2026-06-25).

% ── 1.3  Research Questions and Scope ───────────────────────────────────────
% PURPOSE: State the central RQ, sub-questions, and the (anti-creep) scope.
% MOVE: Operational-narrative objectives (primary -> secondary -> constraint).
% ─────────────────────────────────────────────────────────────────────────

\section{Research Questions and Scope}

The central research question is whether proactive forgetting and consolidation policies
can match the task-resolution performance of full-memory accumulation at materially lower
cost, and --- if any difference exists --- which \emph{content-aware} structure of
forgetting matters most. The framing throughout is that \emph{not every task needs
memory}: the question is not how to remember more, but when and how much to forget without
paying for it in resolved issues. We instantiate the question on the SWE-Bench-CL
continual-learning benchmark for coding agents~\cite{joshi2025}.

This primary question decomposes into five secondary questions, in descending order of
how load-bearing they are for the thesis:

\begin{itemize}
  \item \textbf{SQ1 (trade-off).} How do the six policies compare on continual-learning
        resolution (the CL-F1 proxy) against cost and memory footprint?
  \item \textbf{SQ2 (mechanism).} If a gap exists, what drives it --- the \emph{volume} of
        memory retained, its \emph{content} (which records are kept), or its
        \emph{compression} into consolidated summaries?
  \item \textbf{SQ3 (predictors).} What features of a stored record predict whether it
        will prove \emph{helpful} versus \emph{harmful} when retrieved?
  \item \textbf{SQ4 (growth).} How does the memory footprint evolve over the course of a
        task sequence under each policy?
  \item \textbf{SQ5 (behavioral).} Does accumulation change how the agent behaves
        step to step --- its tool-call count, its error rate?
\end{itemize}

The scope is deliberately narrow, and the narrowness is the contribution. The experiment
covers the eight official SWE-Bench-CL sequences, three seeds, same-repository retrieval
only, and a fixed retrieval depth of $k=5$, on a single frozen model
(\texttt{deepseek-v4-flash}; see the deviations in Section~\ref{sec:deviations}). The
decisive constraint is that \emph{retrieval is held constant across all six policies}:
pure cosine similarity, identical depth and context budget, so that the storage policy is
the only moving part.

The out-of-scope boundaries are conscious limits, not omissions. We do not add conditions
beyond the six pre-registered policies; we do not run a cross-repository main experiment;
and we do not perform a hyperparameter grid search over the policies. These choices keep
the comparison clean enough that its result, whatever it is, can be attributed to
forgetting rather than to an uncontrolled degree of freedom.
